% SPDX-License-Identifier: Apache-2.0
%
% Ethernet in TxHDL. Built by //docs:eth. Every listing is included
% from the tree, and every printout and figure is what the build made
% by running the example.
\documentclass[journal,9pt]{IEEEtran}

% Listings of Rust set by rustfmt at 80 columns: 5.6pt is what fits
% the frame of a column (issue 1061).
\newcommand{\listingsize}{\fontsize{5.6}{6.6}\selectfont}
% SPDX-License-Identifier: Apache-2.0
%
% The house preamble, shared by every document in this directory. Loaded
% after \documentclass, before \title. Keeping it in one file is what
% stops two articles drifting apart in font, listing style or figure
% style.
% Computer Modern. IEEEtran selects Times (ptm) by default, so the face has
% to be chosen explicitly.
%
% Latin Modern is the Computer Modern variety used here, and the choice is
% forced rather than aesthetic. Under T1 encoding, plain `cmr` has no Type 1
% outlines unless cm-super is installed, and it is not in the pinned
% distribution, so pdflatex falls back to EC bitmaps and every glyph in the
% PDF becomes a Type 3 font. That was measured: the first build produced a
% document with 10 Type 3 fonts and no Type 1. Latin Modern is Computer
% Modern redrawn with a full T1 repertoire and real .pfb outlines, so the
% same design comes out scalable.
\usepackage[T1]{fontenc}
\usepackage[utf8]{inputenc}
\usepackage{lmodern}
% microtype is not in the pinned TeX distribution. emergencystretch alone
% keeps the two-column measure from producing overfull lines.
\setlength{\emergencystretch}{3em}

\usepackage{amsmath}
\usepackage{amssymb}
\usepackage{amsthm}
\usepackage{booktabs}
\usepackage{array}
% enumitem is not in the pinned TeX distribution; the list spacing below
% does what \setlist would have done.
\usepackage{float}
\usepackage{xcolor}
\usepackage{listings}
\usepackage{tikz}
\usetikzlibrary{arrows.meta,positioning,fit,backgrounds,calc}
\usepackage[hidelinks,bookmarks=true]{hyperref}

% Numbered, definition-style box for a rule the reader is meant to apply.
\theoremstyle{definition}
\newtheorem{rulebox}{Rule}
\newtheorem{example}{Example}

% Unnumbered asides.
\theoremstyle{remark}
\newtheorem*{tip}{Tip}
\newtheorem*{pitfall}{Pitfall}

% Tighter lists, in place of enumitem's \setlist.
\setlength{\topsep}{2pt}
\setlength{\itemsep}{1pt}
\setlength{\parsep}{0pt}

% Code in the running text. The typewriter face under T1 ligates `<<`
% and `>>` into guillemets, so a shift printed as \code{<<} came out as
% a French quotation mark (issue 571). microtype's \DisableLigatures is
% not in the pinned distribution, but the primitive it rests on is
% pdfTeX's own: \pdfnoligatures strips every ligature from the font it
% is given, and the current font inside \texttt is the typewriter face
% at the size in use, so each size is stripped the first time \code
% sets it. Code wants no ligature at all: two quotes are two quotes.
\newcommand{\code}[1]{\texttt{\ifdefined\pdfnoligatures\pdfnoligatures\font\fi #1}}
\newcommand{\term}[1]{\emph{#1}}

% Syntax highlighting. Four roles, distinguishable in colour and still
% distinguishable in greyscale, because an IEEE article gets printed: the
% keyword is the only bold role, the comment is the only italic one, and the
% two remaining roles differ in lightness as well as in hue.
\definecolor{kwblue}{RGB}{20,60,150}
\definecolor{cmtgreen}{RGB}{90,120,90}
\definecolor{strred}{RGB}{150,50,40}
\definecolor{numgray}{gray}{0.45}
\definecolor{framegray}{gray}{0.70}
\definecolor{bgshade}{gray}{0.97}
% A darker fill for figure cells. The listing background has to stay very
% light to keep code readable; a figure cell has to be clearly shaded or the
% distinction it draws is invisible in print.
\definecolor{cellshade}{gray}{0.90}

% The listing size is the document's choice, because it depends on the
% column: a two-column article sets `\listingsize` to 6pt and keeps its
% listings within 70 columns, or to 5.6pt when it includes source that
% rustfmt sets at 80, which is what fits a column's frame (issue 1061);
% a one-column document keeps the body size and includes source files
% at 80. `//docs:frame_test` checks every listing line against its frame. Lines are never broken; a line that does not fit
% overflows the frame, which is visible, rather than wrapping, which is
% not.
\providecommand{\listingsize}{\scriptsize}

% `'` is deliberately not a string delimiter: the language uses it for loop
% labels, as Rust does, so treating it as a quote would swallow the rest of
% the line.
\lstdefinestyle{txhdl}{
  basicstyle=\ttfamily\listingsize,
  keywordstyle=\bfseries\color{kwblue},
  commentstyle=\itshape\color{cmtgreen},
  stringstyle=\color{strred},
  numberstyle=\tiny\color{numgray},
  morekeywords={mod,use,pub,const,type,struct,enum,trait,impl,fn,async,let,mut,
    match,if,else,loop,while,for,in,break,continue,return,as,await,
    module,bus,channel,transaction,protocol,pipeline,flow,seq,config,
    port,stage,pipe,instance,connect,bind,tag,domain,blackbox,foreign,
    property,role,signal,handler,true,false,
    sequence,interface,modport,implements,package,import,var,wire,func,proc,
    case,default,next,fork,branch,join,state,fsm},
  morecomment=[l]{//},
  morestring=[b]",
  showstringspaces=false,
  columns=fullflexible,
  keepspaces=true,
  breaklines=false,
  % Framed, so a listing is bounded rather than floating in the column.
  frame=single,
  rulecolor=\color{framegray},
  framesep=3pt,
  backgroundcolor=\color{bgshade},
  xleftmargin=3pt,
  xrightmargin=3pt,
  aboveskip=6pt,
  belowskip=6pt,
  % Every listing is numbered and captioned, so it can be referred to by
  % number. `caption` without `float` numbers the listing and leaves it
  % where it was written, which is what a listing in running prose needs.
  captionpos=b,
  abovecaptionskip=2pt,
  belowcaptionskip=6pt,
}

% Figures drawn as text. Framed the same way, and with the highlighting
% switched off, because a box-and-arrow diagram is not code and half its
% words turning blue reads as an accident. Setting a style adds keys rather
% than clearing the ones already set, so the three roles are neutralised by
% hand rather than by leaving them out. Program output is printed in this
% style too, and a netlist's assignment can run past the frame, so a line
% that does not fit is broken and the rest indented; source never is.
\lstdefinestyle{ascii}{
  basicstyle=\ttfamily\listingsize,
  keywordstyle=\ttfamily,
  commentstyle=\ttfamily,
  stringstyle=\ttfamily,
  columns=fullflexible,
  keepspaces=true,
  breaklines=true,
  breakindent=2em,
  frame=single,
  rulecolor=\color{framegray},
  framesep=3pt,
  backgroundcolor=\color{bgshade},
  xleftmargin=3pt,
  xrightmargin=3pt,
  aboveskip=6pt,
  belowskip=6pt,
}
\lstset{style=txhdl}

% House TikZ styles. `shorten` lives on the arrow styles rather than on each
% arrow, so every arrow in the document gets the gap, including ones added
% later. TikZ clips a path at a node's border, so without it an arrowhead
% butts against the box with no gap at all. Keep `node distance` well above
% twice the shorten, or the arrow shrinks to nothing.
\tikzset{
  box/.style={draw,rounded corners=2pt,align=center,inner sep=4pt,
              font=\footnotesize},
  boxf/.style={box,fill=bgshade},
  lbl/.style={font=\scriptsize\itshape,align=center},
  fl/.style={-{Stealth[length=4pt]},shorten >=2.5pt,shorten <=2.5pt},
  fld/.style={fl,dashed},
}


% The contents list: the class gives a section's number 2.75em, which
% a Roman numeral past thirty overruns onto the title, so the box is
% wider here, and a subsection's entry starts where a title does.
\makeatletter
\def\l@section#1#2{\addpenalty{\@secpenalty}\addvspace{1.0em plus 1pt}%
    \@tempdima 5.75em \begingroup \parindent \z@ \rightskip \@pnumwidth%
    \parfillskip-\@pnumwidth {\bfseries\leavevmode #1}\hfil\hbox to\@pnumwidth{\hss #2}\par%
    \endgroup}
\def\l@subsection{\@dottedtocline{2}{5.75em}{5em}}
\def\l@subsubsection{\@dottedtocline{3}{10.75em}{5.5em}}
\makeatother

% SPDX-License-Identifier: Apache-2.0
% The boards' memory maps, written by //tools/memmap from the maps the
% designs decode (issue 444). A document asks for a table with
% \mmget{<what>}{<board>}, and a table the generator does not write
% stops the document rather than leaving a gap.
\input{tools/memmap/mm_tables}
\makeatletter
\newcommand{\mmget}[2]{%
  \@ifundefined{mm@#1@#2}%
    {\PackageError{memmap}{//tools/memmap writes no #1 for #2}%
      {Add it to tools/memmap/main.rs.}}%
    {\csname mm@#1@#2\endcsname}}
\makeatother


\InputIfFileExists{buildstamp.tex}{}{}
\providecommand{\buildstamp}{Draft build (outside the Bazel flow).}

% A range of the Ethernet part, by its begin/end markers.
\newcommand{\ethpart}[3]{%
\lstinputlisting[rangebeginprefix=//\ begin\{,rangebeginsuffix=\},%
rangeendprefix=//\ end\{,rangeendsuffix=\},includerangemarker=false,%
linerange=#1-#1,caption={#2},label={#3}]{lib/parts/src/eth.rs}}

\title{Ethernet in TxHDL}

\author{Filip Filmar and Dragiša Janković%
\thanks{Both authors are independent. Filip Filmar is the
corresponding author, at f@hdlfactory.com; Dragiša Janković is
reached at linkedin.com/in/dragisa-jankovic.}%
\thanks{This is the tutorial on the Ethernet part, for a reader who
knows the AXI link and its AXI-Lite bridge~\cite{axi}. Every listing is
included from the project repository \code{HDL/txhdl} at build time,
and every printout and figure is what the build produced by running
the example. The cover~\cite{cover} lists every document.}%
\thanks{The authors used a large language model, Claude, as an
assistant in exploring the concepts and in writing and constructing
the documents and the programs; every commit of the repository says
so and preserves the prompts that led to it, verbatim.}%
\thanks{\buildstamp}}

\markboth{Ethernet in TxHDL}{Filmar and Janković: Ethernet in TxHDL}

\begin{document}
\maketitle

\begin{abstract}
An Ethernet media access controller (MAC) serializes frames into byte
streams on a physical wire and reconstructs frames from incoming streams.
The TxHDL implementation comprises two core units: a transmitter and a
receiver, each expressed in roughly eighty lines within the lowered
subset. Both units communicate across Gigabit Media Independent Interface
(GMII) links at one byte per clock cycle. Because each unit operates
within a single clock domain, a board implementation clocks each half
directly from its respective physical-layer (PHY) interface. Each unit
buffers a complete frame before transmission or forwarding, decoupling
client channel throughput and facilitating clock-domain crossing. An
accompanying three-word AXI-Lite peripheral bridges transactions from a
processor core or bus host. The build suite tests all three units across a
simulated loopback wire with injected frame corruption, verifying generated
netlists under both \code{nvc} and Verilator. For the Alinx AX7A200B
development board, a companion design integrates the PHY Reduced GMII
(RGMII) interface, clock generation, and clock-domain crossing around the
two MAC halves. Vivado synthesizes, places, and routes this design to a
verified bitstream that successfully echoes traffic on physical hardware.
\end{abstract}

\section{What Is Here}
\label{sec:e-what}

The \code{eth} module in \code{txhdl\_parts} provides three primary units.

\code{EthTx} receives frame bytes across an input channel and transmits them
over GMII using byte bus \code{txd} and transmit enable \code{tx\_en}.
\code{EthRx} samples incoming GMII signals (\code{rxd}, \code{rx\_dv}, and
\code{rx\_er}) and presents validated bytes on an output channel.
\code{EthLite} serves as the register interface: it forwards bytes from
AXI-Lite writes to the transmitter and services AXI-Lite reads with bytes
supplied by the receiver.

Every transaction on either channel conveys an \code{EthByte}, packaging
the data octet alongside an end-of-frame delimiter flag. At this interface,
a frame contains client-visible fields: destination address, source
address, EtherType, and payload data. Physical framing details, including
preamble octets, minimum-length padding, and cyclic redundancy checks,
remain internal to the MAC.

\ethpart{byte}{A byte of a frame, between the MAC and its
client.}{lst:e-byte}

\section{The Frame on the Wire}
\label{sec:e-frame}

On the physical wire, an Ethernet frame consists of seven preamble octets
(\code{0x55}), a start-of-frame delimiter (\code{0xd5}), the MAC frame,
and a four-byte frame check sequence (FCS). Frames shorter than sixty octets
are zero-padded to sixty octets before CRC calculation. An interpacket gap
of twelve idle cycles separates consecutive frames.

The check sequence employs standard CRC-32 arithmetic. The shift register
initializes to \code{0xffffffff}, ingests data least-significant bit first,
and inverts its final state before transmitting the four octets in
little-endian order. When a receiver processes an uncorrupted frame together
with its trailing check sequence, the register converges to the fixed
residue \code{0xdebb20e3}, removing any need for separate comparison
hardware.

Each step right-shifts the register and applies the reversed polynomial
\code{0xedb88320} whenever the expelled bit is set. Processing a full octet
entails XORing the byte into the register and executing eight sequential
steps. Listing~\ref{lst:e-crc} expresses this structure directly through a
single-step primitive iterated eight times.

Earlier compiler passes lacked this concise representation. Lowering
originally inlined function calls by textual expansion at each reference
site. Because each bit step reads its argument twice, the eighth iterated
call replicated the initial register expression $2^8 = 256$ times, causing
multi-minute compile delays. Hardware designers previously worked around
this expansion by computing parallel 32-bit linear parity equations
directly, obscuring the underlying shift algorithm. With common
subexpression elimination now allocating intermediate wires for
multiply-referenced expressions (issue~126), the code retains its iterative
form while adding only sixteen internal wires to the netlist. A reference
test verifies this logic against bitwise software evaluation across four
register patterns.

\ethpart{crc}{A byte into the CRC-32 register: one step, and eight of
them.}{lst:e-crc}

\section{The Transmitter}
\label{sec:e-tx}

Because physical Ethernet transmission cannot pause mid-frame, the
transmitter buffers each complete frame before initiating wire output.
While idle and awaiting data, it ingests available bytes from channel
\code{tx} into a 2048-byte buffer named \code{frame}. Receiving an octet
with its terminal flag asserted transitions status signal \code{whole}
high.

Asserting \code{whole} triggers the output state machine: emitting
preamble octets, serializing payload bytes while updating CRC state,
padding to sixty octets if necessary, transmitting the CRC-32 value, and
observing the required interpacket gap. Upon concluding the gap, the
transmitter asserts \code{sent}, clears buffer occupancy, and accepts
subsequent inbound traffic.

The transmitter architecture partitions into two concurrent processes.
The ingestion process consumes one byte per clock cycle through a
single-wait loop. The wire output process expresses the transmission
sequence directly: waiting for a complete frame, generating the preamble
sequence, iterating over buffered data up to index \code{fill}, streaming
zero padding for short frames, dispatching the four CRC octets, and
counting interpacket gap cycles. Lowering synthesizes these sequential
waits into an explicit state machine and allocates loop counter registers,
eliminating the manually managed \code{phase} and \code{pos} registers
used in early revisions.

Clients may supply payload bytes at arbitrary rates, with backpressure
asserted whenever an active transmission occupies the buffer. This rate
decoupling permits asynchronous clock-domain crossing between peripheral
logic and the physical transmit clock on target hardware.

\ethpart{txstate}{The transmitter's state.}{lst:e-txstate}

\ethpart{tx}{The transmitter's step.}{lst:e-tx}

\section{The Receiver}
\label{sec:e-rx}

Because physical wire traffic cannot pause for backpressure, the receiver
buffers an entire frame before presenting validated data to downstream
consumers. The input monitor detects data valid signal \code{rx\_dv},
scans past preamble bytes until recognizing the start-of-frame delimiter on
\code{rxd}, and buffers one byte per cycle while \code{rx\_dv} remains
asserted. Throughout reception, the CRC register computes over all
incoming octets including the trailing checksum. When \code{rx\_dv}
negates, the receiver validates the frame: the CRC register must match the
expected residue, error flag \code{rx\_er} must remain unasserted, and
total frame length must exceed four bytes. Validated frames are then
presented across channel \code{rx} stripped of their check sequence, with
the final payload byte tagged accordingly.

Counter \code{dropped} records discarded frames across three distinct
error conditions. First, frames failing the CRC residue check are dropped
immediately. Second, frames arriving while previous data remains pending
on the client channel are discarded due to buffer occupancy. Third,
incomplete frame fragments encountered without a valid preamble and start
delimiter are ignored until \code{rx\_dv} deasserts.

The receiver likewise separates into two concurrent processes. The
physical interface process monitors wire signals on every cycle: hunting
for frame boundaries, streaming bytes into memory, evaluating integrity,
and asserting signal \code{full}. The client drain process executes
sequentially: awaiting \code{full}, driving the payload size onto
\code{rx\_len}, streaming payload octets bounded by \code{len - 4} until
consumed, and asserting \code{given} to release the buffer for subsequent
traffic. As with the transmitter, the lowering generates dedicated state
registers and loop counters automatically.

\ethpart{rxstate}{The receiver's state.}{lst:e-rxstate}

\ethpart{rx}{The receiver's step.}{lst:e-rx}

\section{The Peripheral}
\label{sec:e-lite}

\code{EthLite} is an AXI-Lite peripheral of three words, moving one byte per
transaction. Its words are declared by the \code{regmap!} macro at the head
of Listing~\ref{lst:e-lite}: read-only register \code{status}, indicating
buffered receive data and transmit availability; write-only register
\code{txbyte}, supplying payload data and end-of-frame flags; and
register \code{rxbyte}, whose read pops queued octets alongside status and
validity bits. This macro definition generates read multiplexing logic,
address decode strobes, and host-accessible register offsets from a single
specification.

AXI-Lite write responses complete only after the transmitter accepts the
queued byte, stalling fast hosts via bus backpressure to prevent buffer
overruns. Reading the receive register extracts the oldest queued octet
with status bit~9 asserted; if the FIFO is empty, the read completes with
bit~9 cleared without altering state. An interrupt request line asserts
whenever received bytes remain available.

Single-byte transactions introduce overhead on 1500-byte Ethernet payloads,
but bridge optimizations mitigate transaction latency. When addressed with
fixed-burst transfers, the AXI-to-AXI-Lite bridge converts each beat into
back-to-back single-cycle operations, allowing host software to stream
complete frames to \code{txbyte} or burst 16-byte chunks from
\code{rxbyte} efficiently.

\ethpart{lite}{The peripheral.}{lst:e-lite}

\section{The Run}
\label{sec:e-run}

Example \code{ex\_eth} instantiates the three units behind
\code{LiteBridge<1, ..>} and a host transaction tracker, looping the
transmitter GMII interface directly into receiver inputs with a one-cycle
delay. The simulated host issues three frames as fixed-burst transfers.
Following each frame, the host reads 16-byte bursts from \code{rxbyte} until
the frame returns or twelve successive bursts return empty. Simulated
loopback wiring inverts one data byte within the second frame during
transmission to test error handling.

The first frame spans 20 payload bytes and returns padded to the 60-byte
minimum. The corrupted second frame fails its check sequence and is
rejected, incrementing the receiver's drop counter. The third frame
contains 72 bytes and returns unmodified. The run asserts all three
transmission outcomes, confirming that three frames were sent and exactly
one was dropped.

\begin{figure*}[t]
\centering
\input{eth_timing.tex}
\caption{The first frame: its bytes offered to the transmitter, the
frame whole and then sent, \code{tx\_en}, the receiver's
\code{rx\_dv}, the frame being received and then full, and the bytes
offered back, with the interrupt line.}
\label{fig:e-run}
\end{figure*}

\lstinputlisting[caption={\code{ex\_eth.rs}. Run with
\code{bazel run //lib/examples:ex\_eth}.},%
label={lst:e-ex}]{lib/examples/ex_eth.rs}

\lstinputlisting[style=ascii,caption={What \code{ex\_eth} printed when
the build ran it.},label={lst:e-out}]{out_eth.txt}

The transmitter, the receiver, and the peripheral are lowered, and the build
simulates each netlist against this run under \code{nvc} and under Verilator
at its boundary ports. Unit tests within the module verify transmit wire
formatting across frame lengths of 59, 60, and 61 bytes, confirming correct
behavior under each drop condition.

\section{On the Board}
\label{sec:e-board}

The Alinx AX7A200B provides a single gigabit Ethernet port wired to a
JL2121-N040I PHY over RGMII. The physical layer transceiver negotiates
10, 100, or 1000~Mbit operation and duplex modes autonomously based on
hardware strap pins sampled during power-up. While documentation indicates
default management bus responsiveness at MDIO address~3, board probing
identified the PHY at address~0. RGMII transfers GMII octets as dual-edge
nibbles: the lower nibble on the rising clock edge and the upper nibble
on the falling clock edge. Similarly, the control signal encodes data
valid on rising edges and an XOR error condition on falling edges.
Handling dual-data-rate signaling requires vendor-specific FPGA primitives
(\code{ODDR} and \code{IDDR}) not directly generated by the lowered
subset, prompting a handwritten Verilog shim in \code{//eth/board}. The
board-level demonstration implements an echo pipeline, immediately
retransmitting each uncorrupted received frame.

On the flagship system, the core accesses the port through \code{EthSlots},
the slot register interface expected by the Zephyr driver, located at
peripheral offset \mmget{at}{ethslots}. The register peripheral
\code{EthLite} described above is the interface used by \code{ex\_eth}.
The flagship memory map details this assignment.

Module \code{eth\_rgmii.v} implements the RGMII wrapper. Transmit clock
generation routes through an \code{ODDR} primitive identical to data lines,
keeping output edges aligned and allowing internal PHY delay lines to
center the clock within data eyes. Conversely, the incoming receive clock
passes through an MMCM configured for a 90-degree phase shift (2~ns of the
8~ns period), placing the sampling edge safely inside the valid data
window before driving the receive MAC logic. Both phase alignments were
determined empirically during board bring-up (issue~231).

\lstinputlisting[caption={\code{eth/board/eth\_rgmii.v}: RGMII to
GMII.},label={lst:e-rgmii}]{eth/board/eth_rgmii.v}

Because the receiver and transmitter operate in separate clock domains
(PHY clock and internal FPGA clock), data movement between them crosses
domains through an asynchronous FIFO. Module \code{chan\_cdc.v} in
\code{//lib/board} provides this channel: a sixteen-entry FIFO using
standard valid/ready handshake signals on each side. Gray-code pointers
cross domains through two flip-flop synchronizers. Buffering complete
frames allows the FIFO to hold bytes for several cycles without data loss.

Each MAC domain receives an asynchronous reset synchronized to its local
clock: the clock manager lock signal passes through dual flip-flops in each
domain, aligning release with local clock edges. Because the PHY clock
remains inactive until link startup, its synchronizer maintains active reset
from power-on until clock edges arrive. Lowered modules reset strictly
through their \code{rst} input (issue~509). The cross-domain FIFO samples
both resets, remaining empty until both domains release cleanly
(issue~1325).

\lstinputlisting[caption={\code{lib/board/chan\_cdc.v}: a channel from
one clock to another.},label={lst:e-cdc}]{lib/board/chan_cdc.v}

\lstinputlisting[caption={\code{eth/board/eth\_echo.v}: the echo
design.},label={lst:e-echo}]{eth/board/eth_echo.v}

The command \code{bazel build //eth:echo\_synth} runs Vivado synthesis
for the target device without errors or critical warnings. Vivado emits
only two benign warnings: unused frame counters are trimmed, and
unconnected management pins are tied off.

Next, \code{bazel build //eth:echo\_pnr} places and routes the design
against board constraints to generate the bitstream. All timing constraints
close cleanly, showing 1.491~ns setup slack and 0.065~ns hold slack across
6549 endpoints with clean design-rule checks. The design consumes 1330
LUTs (776 for frame and FIFO memories), 232 flip-flops, and 19 pins.
Because external board trace delays are unconstrained here, reported slack
applies strictly to internal FPGA paths.

Listing~\ref{lst:e-xdc} gives the clock and LED constraints in
\code{eth/board/ax7a200\_eth.xdc}. Listing~\ref{lst:e-pins} supplies the
fifteen PHY pins from the AX7A200B manual in \code{eth/board/eth\_pins.xdc},
which the flagship system also imports.

\lstinputlisting[caption={\code{eth/board/ax7a200\_eth.xdc}: the
board's pins.},label={lst:e-xdc}]{eth/board/ax7a200_eth.xdc}

\lstinputlisting[caption={\code{eth/board/eth\_pins.xdc}: the PHY's
pins.},label={lst:e-pins}]{eth/board/eth_pins.xdc}

\section{What Is Not Done}
\label{sec:e-not}

The echo design was validated on physical hardware. On September 19,
2026, the AX7A200B established a gigabit link with a host and reflected
1000 test frames: 500 of 1400 bytes and 500 of 4 bytes. Every frame
returned intact with short frames padded to 60 bytes, and host network
diagnostics counted zero frame check sequence errors.

Achieving clean communication required two pin timing corrections
(issue~231). The receive path originally clocked input pins through an
unshifted global buffer, preventing data lock. Passing the receive clock
through an MMCM introduces a dedicated phase shift: 90~degrees (2~ns of the
8~ns period) in the echo design and 78.75~degrees in the flagship.
Conversely, the transmit path initially delayed the clock by a quarter
cycle. Because the PHY adds internal delay, this combined shift aligned
clock edges with data transitions, causing bit error rates near
$10^{-4}$ under \code{ethtool -S rx\_errors}. Realigning transmit clock
and data edges at FPGA outputs eliminated all errors, matching vendor
reference designs and yielding loss-free ping traffic.

The receive phase requirement is understood partially (issue~864).
Observed timing fits a PHY that centers its receive clock within data eyes,
placing clock edges roughly 2~ns after data transitions. The PHY answers
MDIO at address~0 with identifier \code{937c 4032}, identifying the JLSemi
gigabit family. Linux driver sources place receive delay configuration at
bit~9 of register~17 on page~3336. Because reading this location requires
writing page registers and data sheets are unavailable, power-on default
register values remain unconfirmed.

Both devices use 3.3~V LVCMOS signaling: core board banks and carrier
PHY pins are tied to 3.3~V rails. Test target \code{//eth/hil:run}
executed these hardware verification runs over physical cabling.

The current MAC operates only at gigabit line rate. At 10 or 100~Mbit, the
PHY transmits nibbles across multiple cycles, which this single-cycle byte
pipeline does not process. Because the autonomous PHY negotiates all three
speeds, link partners must be configured for gigabit operation. The echo
design leaves management pins idle, relying on power-on strap defaults
without active link monitoring. The flagship SoC provides an MDIO master
at peripheral offset \code{0x3700} (issue~864).

Each MAC half buffers a single frame. If client logic fails to drain
the receiver before another packet arrives, the arriving frame is dropped.
Peripheral \code{EthLite} transfers one byte per bus transaction. The
Vreteno system bypasses byte-by-byte register transfers, accessing MAC
ports through \code{EthSlots} in slot~5 via DMA engines that stream
packets directly to and from DDR3 memory (issue~151).

\begin{table}[htbp]
\caption{Port performance on hardware across DMA engines and CPU copy routines
for 1000 frames (2026-10-04). Transmit was unconstrained; receive was paced
at 25\,020 and 16\,681 frames per second.}
\label{tab:e-perf}
\centering
\begin{tabular}{@{}llrrr@{}}
\toprule
Mechanism & Bytes & Frames/s & MB/s & Cycles \\
\midrule
send, engine & 60 & 90\,001 & 5.4 & 155 \\
send, copy & 60 & 20\,787 & 1.2 & 3\,735 \\
send, engine & 1514 & 8\,358 & 12.7 & 158 \\
send, copy & 1514 & 1\,104 & 1.7 & 84\,663 \\
receive, engine & 60 & 24\,989 & 1.5 & 883 \\
receive, copy & 60 & 12\,545 & 0.8 & 4\,247 \\
receive, engine & 1514 & 16\,672 & 25.2 & 884 \\
receive, copy & 1514 & 913 & 1.4 & 105\,059 \\
\bottomrule
\end{tabular}
\end{table}

Table~\ref{tab:e-perf} compares DMA engine performance against programmed CPU
copy operations measured by \code{ethperf.rs} on the flagship build
(issues 1022 and 1038). CPU copies establish an empirical baseline; single-byte
\code{EthLite} bus accesses would incur even higher overhead. The final
column reports CPU active cycles per frame, excluding idle wait time.

With DMA engines, CPU overhead remains constant across packet sizes: roughly
155 cycles to dispatch a transmit frame and 880 cycles to retire a received
packet. Programmed copying requires 225 cycles per 32-bit DDR3 word across the
bus bridge, consuming 85\,000 cycles to send and 105\,000 cycles to receive a
full frame. Transmitting maximum-size frames via DMA sustains 12.7~MB/s,
limited by DDR3 memory read latency (issue~1023). Inbound DMA achieves
25.2~MB/s with zero packet loss at test pacing, while 60-byte bursts experience
a 0.1\% drop rate. Under unthrottled line-rate traffic, the single-buffer
receiver sustains roughly one full packet in five because wire input cannot
be stalled. Earlier test runs occasionally deposited unannounced packets into
unserviced buffer slots. Following issue~1313, frames arriving when both
slots remain occupied are dropped and counted under \code{rx\_errors}, as
detailed in the \code{EthSlots} section of the examples document. Every
transmitted packet arrived verified at the link partner.

Implementing the Ethernet module resolved three compiler lowering issues.
First, signals matching HDL keywords (\code{next}, \code{out}, \code{begin},
\code{body}, \code{buf}, or \code{byte}) originally produced invalid syntax;
lowering now escapes reserved identifiers systematically as \code{next\_rw}
(issue~497). Second, bit selections within channel fields (\code{wh.data.bit(8)})
emitted nested part-selects rejected by Verilog; lowering now resolves field
bits to absolute vector indices. Third, integer literals above $2^{31}-1$
were emitted as unsized values, overflowing signed 32-bit integers in Verilog
and default ranges in VHDL. Lowering now emits constants with explicit bit
widths, using sized hexadecimal in Verilog and bit strings in VHDL, allowing
direct residue comparison against \code{0xdebb\_20e3}.

\IEEEtriggeratref{1}
\begin{thebibliography}{9}

\bibitem{axi}
F.~Filmar and D.~Janković, ``AXI in TxHDL,'' project repository
\code{HDL/txhdl}, \code{//docs:axi}, 2026.

\bibitem{cover}
F.~Filmar and D.~Janković, ``TxHDL documents,'' project repository
\code{HDL/txhdl}, \code{//docs:cover}, 2026.

\end{thebibliography}

\end{document}
